%%
%% Appendix A: RTFM XD-I3D reproduction investigation
%%
% Source (narrative + all reproduction numbers):
%   .planning/phases/04b-xd-violence-main-results-rtfm-xd-i3d-gate/04b-05-SUMMARY.md (tracked)
%   .planning/phases/07-xd-violence-hyperparameter-sweep/07-VERIFICATION.md (tracked)
% Per-category table: thesis/tables/tab_percat_rtfm.tex (Class R, generated by
%   scripts/thesis_per_category_tables.py from tracked per_category.csv files; s42-only).
\chapter{RTFM XD-I3D Reproduction Investigation}
\label{app:a}

The experimental plan for this thesis included a baseline-reproduction gate:
before trusting our own MIL evaluation harness on XD-Violence, we attempted
to reproduce RTFM~\cite{tian2021rtfm} on XD-Violence using pre-extracted I3D
RGB features. The attempt missed its acceptance gate by a wide margin, the
miss was accepted and disclosed as a limitation
(Chapter~\ref{ch:08}), and a structured investigation traced the gap to
training-regime differences rather than to any defect in the evaluation
pipeline. This appendix documents that investigation in full---the result,
the diagnostics that ruled out the competing explanations, and the reasons
the gap does not affect the dual-modal fusion results reported in
Chapters~\ref{ch:05} and~\ref{ch:06}.

\section{Protocol and Acceptance Gate}
\label{sec:appa-protocol}

The reproduction trained an RTFM-style magnitude-based MIL head on
XD-Violence using a pre-extracted five-crop I3D RGB feature cache
(1{,}024-dimensional snippet features; features are consumed frozen, exactly
as in the rest of this thesis, cf.\ Chapter~\ref{ch:04}). Each training batch
paired 3 normal with 3 abnormal videos, with the five spatial crops treated
as independent bags (an effective 15+15 bags per step). Evaluation used the
same frame-level protocol as our XD-Violence experiments: snippet scores
expanded to frames against the official temporal annotations, 800 test
videos.

The pre-registered acceptance gate was the published RTFM XD-Violence
I3D-RGB result of 77.81\% AP minus a one-point tolerance, i.e.\ a gate
threshold of 0.7681 frame-level AP.

\section{Result: An Accepted Gate Miss}
\label{sec:appa-result}

% Source: 04b-05-SUMMARY.md empirical-results block (tracked).
\begin{table}[t]
\centering
\caption{RTFM XD-I3D-RGB reproduction against the published anchor. The
shortfall is quoted against the pre-registered gate threshold (published
anchor minus one point). Single run, seed 42.}
\label{tab:appa-gate}
\begin{tabular}{lr}
\toprule
Quantity & Value \\
\midrule
Published RTFM XD I3D-RGB anchor (AP) & 0.7781 \\
Pre-registered gate threshold (anchor $-$ 1 point) & 0.7681 \\
Reproduction frame-level AP & \textbf{0.6570} \\
Shortfall vs.\ gate threshold & $-$11.11 pp \\
\midrule
Reproduction frame-level AUC & 0.8675 \\
Reproduction snippet-level AUC & 0.8699 \\
$|$frame AUC $-$ snippet AUC$|$ & 0.0024 \\
Test videos evaluated & 800 / 800 \\
\bottomrule
\end{tabular}
\end{table}

Table~\ref{tab:appa-gate} summarizes the outcome. The reproduction reached a
frame-level AP of \textbf{0.6570} against the published anchor of 0.7781,
an 11.11-point shortfall relative to the gate threshold. The gate was
declared a miss, and the miss was \emph{accepted}: rather than iterating on
the reproduction until the number cooperated, we documented the gap, ran the
diagnostics below to establish its cause, and carried the result forward as
a disclosed limitation.

One competing explanation can be dismissed immediately from
Table~\ref{tab:appa-gate} itself. If the snippet-to-frame score expansion
were misaligned (scores assigned to the wrong frames), frame-level and
snippet-level AUC would diverge; they differ by only 0.0024, far below the
0.02 sanity ceiling used throughout this thesis. The evaluation-side
expansion is correct.

% Source: generated by scripts/thesis_per_category_tables.py (CPU-only, tracked).
% Data: results/xd_i3d_rtfm_i3d_s42/per_category.csv +
%       results/xd_i3d_rtfm_i3d_flow_s42/per_category.csv (tracked).
% SEED 42 ONLY -- the RTFM reproduction and its Flow-swap diagnostic are single-seed runs.
\begin{table}[t]
\centering
\caption{Per-category frame-level AP (\%) on XD-Violence for the RTFM-I3D reproduction (RGB features) and its Flow-swap diagnostic. Seed 42 only (single-seed runs).}
\label{tab:percat-rtfm}
\begin{tabular}{llrr}
\toprule
Code & Category & RTFM (RGB) AP (\%) & Flow-swap AP (\%) \\
\midrule
B1 & Fighting & 63.24 & 53.43 \\
B2 & Shooting & 39.69 & 34.17 \\
B4 & Riot & 71.57 & 61.73 \\
B5 & Abuse & 4.34 & 4.71 \\
B6 & Car Accident & 25.06 & 19.57 \\
G & Explosion & 51.95 & 25.44 \\
\bottomrule
\end{tabular}
\end{table}


Table~\ref{tab:percat-rtfm} breaks the reproduction down by XD-Violence
category (seed 42 only; the reproduction and its diagnostic are single-seed
runs). The profile is informative: Riot (71.57\% AP) and Fighting (63.24\%)
are handled reasonably, while Abuse collapses to 4.34\% AP---despite a
per-category AUC of 0.8025 for the same class---reflecting Abuse's very low
positive-frame prevalence, which frame-level AP punishes much harder than
AUC. The reproduction's weakness is concentrated where positive frames are
rare and events are short, consistent with a training-regime (rather than
evaluation) deficit: the model ranks frames sensibly but does not separate
rare positives sharply enough for precision-based metrics.

\section{Ruling Out Data Coverage: the Flow-Swap Diagnostic}
\label{sec:appa-flow}

Inspection of the RGB feature cache revealed a genuine data problem: the
pre-extracted I3D RGB training cache was truncated at the alphabetical tail
of the corpus (video names beginning V/W/Y), leaving 2{,}748 of 3{,}360
training videos (81.8\%) with features; validation coverage was 477/594
(80.3\%). The test-side cache was complete (800/800), so evaluation was
unaffected, but an 18\% training-data loss was a plausible cause of the gate
miss.

The corresponding optical-flow I3D cache, by contrast, was verified 100\%
complete and is layout-identical (same shapes, same five-crop structure).
This enabled a decisive, zero-code diagnostic: swap the Flow cache in place
of the RGB cache (via a filesystem junction, so the training pipeline is
byte-for-byte the same) and re-run the identical training and evaluation
recipe. If the missing 18\% of training data caused the miss, the
full-coverage Flow run should close much of the gap.

% Source: 04b-05-SUMMARY.md flow-diagnostic block (tracked).
\begin{table}[t]
\centering
\caption{Flow-swap diagnostic: identical training recipe on the truncated
RGB cache vs.\ the complete Flow cache. Full data coverage scores
\emph{worse}, ruling out the cache truncation as the cause of the gate miss.
Seed 42.}
\label{tab:appa-flow}
\begin{tabular}{lrrr}
\toprule
Metric & RGB (81.8\% train) & Flow (100\% train) & $\Delta$ \\
\midrule
Frame-level AP & 0.6570 & \textbf{0.5916} & $-$6.54 pp \\
Frame-level AUC & 0.8675 & 0.8341 & $-$3.33 pp \\
Snippet-level AUC & 0.8699 & 0.8365 & $-$3.34 pp \\
Train coverage & 2748/3360 & 3360/3360 & --- \\
Test coverage & 800/800 & 800/800 & --- \\
\bottomrule
\end{tabular}
\end{table}

Table~\ref{tab:appa-flow} shows the opposite: with 100\% training coverage,
the Flow run scored 0.5916 AP---6.54 points \emph{worse} than the truncated
RGB run. (The Flow run's per-category profile appears alongside the RGB
profile in Table~\ref{tab:percat-rtfm}.) Data completeness is therefore not
the bottleneck. The diagnostic also disqualified the most tempting remedy:
concatenating RGB and Flow features into a 2{,}048-dimensional input, the
standard RTFM configuration, is unlikely to rescue the reproduction when the
Flow stream alone is substantially weaker than the truncated RGB stream in
this pipeline.

\section{Root-Cause Diagnostics}
\label{sec:appa-rootcause}

% Source: 07-VERIFICATION.md (tracked) — Phase-7 diagnostic findings:
% temporal-interpolation check PASS, I3D feature audit (1024-d vs 2048-d),
% published-code comparison (3 high-impact training-regime differences).
A second round of diagnostics, run during the later hyperparameter-sweep
phase, examined the reproduction pipeline against the published RTFM
implementation directly. Three findings jointly explain the gap:

\begin{enumerate}
    \item \textbf{Feature mismatch.} Our pre-extracted cache provides
    1{,}024-dimensional I3D RGB features, whereas the published RTFM
    implementation consumes 2{,}048-dimensional I3D features. The cached
    features were produced by a different extraction pipeline and are not
    byte-compatible with the features the published numbers were obtained
    on; frozen-feature reproductions inherit every upstream difference in
    the extractor.
    \item \textbf{Simplified model variant.} Our reproduction implements the
    magnitude-based top-$k$ MIL objective but omits RTFM's multi-scale
    temporal network (MTN), a component the original ablations credit with
    a substantial share of its performance. The simplification was
    deliberate---the reproduction's purpose was to sanity-check the MIL
    evaluation stack, not to re-engineer RTFM---but it is a real capacity
    difference.
    \item \textbf{Training-regime differences.} The reproduction trains
    with 3 video pairs per batch (15+15 bags after crop-as-sample
    expansion) versus RTFM's batch size of 32, along with the associated
    differences in optimization dynamics.
\end{enumerate}

A separate temporal-interpolation check on the evaluation side passed: the
snippet-to-frame alignment used to score XD-Violence is correct. Combined
with the snippet-vs-frame AUC agreement of Section~\ref{sec:appa-result},
the diagnostics localize the entire gap on the training side. The
investigation's one-line conclusion, carried into the thesis record
verbatim, is: \emph{training regime, not eval bug}.

\section{Disposition and Why Fusion Results Are Unaffected}
\label{sec:appa-disposition}

The pre-registered fallback cascade for a gate miss was evaluated
explicitly. Relaxing the anchor was inapplicable---the policy permitted
relaxation only for misses under 3 points, and the shortfall is 11.11 points
against the gate; even against the corrected secondary anchor (MGFN on the
same I3D features, 79.19\% AP), the gap is 13.49 points. RGB+Flow
concatenation was rejected on the strength of the flow-swap evidence.
Adding the MTN module was rejected on budget grounds: an estimated 2--3 days
of engineering better spent on the TTA investigation that constitutes this
thesis's original contribution. The remaining option---accept and document
as a limitation---was selected.

For calibration, published reproduction attempts of RTFM on XD-Violence
land below the paper's number with some regularity: the VadCLIP authors
reported their own RTFM reproduction at roughly 74--76\% AP, and public
issue threads on the official repository report reproductions in the
72--77\% range. Our 65.70\% sits below that band, which is consistent with
the compound mismatch documented above (different feature extraction,
missing MTN, smaller batches) rather than with a single benign discrepancy.

Two consequences matter for the rest of this thesis. First, the fusion
results of Chapters~\ref{ch:05} and~\ref{ch:06} are unaffected: they share
only the evaluation harness with this reproduction, and that harness is
exactly the component the diagnostics cleared (snippet-to-frame expansion
and temporal alignment both verified). Second, wherever RTFM appears in the
comparison tables of Chapter~\ref{ch:07}, we cite the \emph{published}
RTFM numbers, never our reproduction---the reproduction is reported here
solely as a documented negative result on reproducing a third-party
baseline under a mismatched feature pipeline.
